\documentclass[twoside,a4paper,12pt]{article}

\usepackage{EthischeFragenBeimAutomatisiertenFahren}

\begin{document}

\thispagestyle{empty}
\begin{center}
	\includegraphics[scale=1]{resources/fernunisignet-sw}
\end{center}
\vskip 4cm
\begin{center}
	{\Large Seminararbeit 1902 WS 2018/19\par}
	\vskip 2.5cm
	{\textbf{\Large Ethische Fragen} \par}    
	\vskip 0.2cm	
	{\textbf{\Large beim automatisierten Fahren} \par}    
	\vskip 1cm
	{\Large Nikolai Kabiolskij\\Christian Oevermann \par}
	\vskip 4.5 cm
	\Large Dezember 2018
\end{center}
\vspace*{\stretch{1}}
%%\FUiH\\ ich kann damit nicht kompilieren und weiss nicht was das bewirken muss, sonst hätte ich es irgendwie ersetzt
Fakultät für Mathematik und Informatik\\
Lehrgebiet Kooperative Systeme\\
Betreut von apl. Prof. Dr. Christian Icking und Dr. Lihong Ma\\
  
\newpage
\shipout\hbox{}
\newpage

\frontmatter

\section*{Einleitung}
Seit je und immer wollten die Menschen sich fortbewegen können und haben Ihre Mobilität zunehmend ausgebaut. Mit dem rasanten Anstieg von Technologien und der zunehmenden Digitalisierung in der Gesellschaft entstehen neue Möglichkeiten, welche man für die Mobilität der Zukunft einsetzen kann und soll. Das Thema dieser Arbeit ist das autonome Fahren, im Konkreten werden die ethischen Fragen, welche dabei entstehen erläutert. \\ Auch wenn das autonome Fahren als sicher zu erscheinen vermag, da hierfür nur sichere Systeme eingesetzt werden sollen, ist die Möglichkeit eines Unfalls dennoch nicht ausgeschlossen. Es können ebenso verschiedene Dilemma Situationen auftreten bei denen keine der möglichen Entscheidungen hundertprozentig richtig ist. Ebenso ist die Richtigkeit der Entscheidung von Mensch zu Mensch unterschiedlich und verleiht uns allen eine besondere Persönlichkeit. \\ Um all die ethischen Fragen rund um autonomes Fahren zu beantworten oder zumindest anzugehen wurde von dem damaligen Bundesverkehrsminister Herr Alexander Dobrindt eine Ethik-Kommission zum automatisierten Fahren eingesetzt. Geleitet wurde die Kommission von dem ehemaligen Bundesverfassungsrichter Prof. Dr. Dr. Udo Di Fabio. Die Ethik-Kommission des BMVI setzte sich aus 14 Wissenschaftlern und Experten aus den Fachrichtungen Ethik, Recht und Technik zusammen. \cite{site1} Ziel der Kommission war, die Leitlinien für die Entwicklung von autonom fahrenden Systemen zu erarbeiten. \\
Das Ergebnis der fünf Sitzungen der Kommission sind die 20 Leitlinien, welche in dieser Arbeit näher betrachtet werden. 

\newpage

\tableofcontents

\newpage

\mainmatter

\cleardoublepage
\section{Hauptteil}

\begin{quote}
\glqq
1. Teil- und vollautomatisierte Verkehrssysteme dienen zuerst der Verbesserung der Sicherheit aller Beteiligten im Straßenverkehr. 
Daneben geht es um die Steigerung von Mobilitätschancen und die Ermöglichung weiterer Vorteile. Die technische Entwicklung 
gehorcht dem Prinzip der Privatautonomie im Sinne eigenverantwortlicher Handlungsfreiheit.\grqq~\cite{ek}
\end{quote}

\begin{quote}
\glqq
2. Der Schutz von Menschen hat Vorrang vor allen anderen Nützlichkeitserwägungen. Ziel
ist die Verringerung von Schäden bis hin zur vollständigen Vermeidung. Die Zulassung
von automatisierten Systemen ist nur vertretbar, wenn sie im Vergleich zu menschlichen
Fahrleistungen zumindest eine Verminderung von Schäden im Sinne einer positiven Risikobilanz verspricht.\grqq~\cite{ek}
\end{quote}

\begin{quote}
\glqq
3. Die Gewährleistungsverantwortung für die Einführung und Zulassung automatisierter
und vernetzter Systeme im öffentlichen Verkehrsraum obliegt der öffentlichen Hand.
Fahrsysteme bedürfen deshalb der behördlichen Zulassung und Kontrolle. Die Vermeidung von Unfällen ist Leitbild, wobei technisch 
unvermeidbare Restrisiken einer Einführung des automatisierten Fahrens bei Vorliegen einer grundsätzlich positiven Risikobilanz
nicht entgegenstehen.\grqq~\cite{ek}
\end{quote}

\begin{quote}
\glqq
4. Die eigenverantwortliche Entscheidung des Menschen ist Ausdruck einer Gesellschaft, in
der der einzelne Mensch mit seinem Entfaltungsanspruch und seiner Schutzbedürftigkeit
im Zentrum steht. Jede staatliche und politische Ordnungsentscheidung dient deshalb
der freien Entfaltung und dem Schutz des Menschen. In einer freien Gesellschaft erfolgt
die gesetzliche Gestaltung von Technik so, dass ein Maximum persönlicher Entscheidungsfreiheit in einer allgemeinen 
Entfaltungsordnung mit der Freiheit anderer und ihrer
Sicherheit zum Ausgleich gelangt.\grqq~\cite{ek}
\end{quote}

\begin{quote}
\glqq
5. Die automatisierte und vernetzte Technik sollte Unfälle so gut wie praktisch möglich vermeiden. Die Technik muss nach 
ihrem jeweiligen Stand so ausgelegt sein, dass kritische
Situationen gar nicht erst entstehen, dazu gehören auch Dilemma-Situationen, also eine
Lage, in der ein automatisiertes Fahrzeug vor der \glqq Entscheidung\grqq\ steht, eines von zwei
nicht abwägungsfähigen Übeln notwendig verwirklichen zu müssen. Dabei sollte das gesamte Spektrum technischer 
Möglichkeiten --- etwa von der Einschränkung des Anwendungsbereichs auf kontrollierbare Verkehrsumgebungen, 
Fahrzeugsensorik und Bremsleistungen, Signale für gefährdete Personen bis hin zu einer Gefahrenprävention mittels
einer \glqq intelligenten\grqq\ Straßen-Infrastruktur --- genutzt und kontinuierlich weiterentwickelt
werden. Die erhebliche Steigerung der Verkehrssicherheit ist Entwicklungs- und Regulierungsziel, und zwar bereits in der 
Auslegung und Programmierung der Fahrzeuge zu defensivem und vorausschauendem, schwächere Verkehrsteilnehmer (\glqq Vulnerable Road
Users\grqq) schonendem Fahren.\grqq~\cite{ek}
\end{quote}

\begin{quote}
\glqq
6. Die Einführung höherer automatisierter Fahrsysteme insbesondere mit der Möglichkeit
automatisierter Kollisionsvermeidung kann gesellschaftlich und ethisch geboten sein,
wenn damit vorhandene Potentiale der Schadensminderung genutzt werden können.
Umgekehrt ist eine gesetzlich auferlegte Pflicht zur Nutzung vollautomatisierter Verkehrssysteme oder die Herbeiführung 
einer praktischen Unentrinnbarkeit ethisch bedenklich, wenn damit die Unterwerfung unter technische Imperative verbunden 
ist (Verbot der Degradierung des Subjekts zum bloßen Netzwerkelement).\grqq~\cite{ek}
\end{quote}

\begin{quote}
\glqq
7. In Gefahrensituationen, die sich bei aller technischen Vorsorge als unvermeidbar erweisen, besitzt der Schutz menschlichen 
Lebens in einer Rechtsgterabwägung höchste Priorität. Die Programmierung ist deshalb im Rahmen des technisch Machbaren 
so anzulegen, im Konflikt Tier- oder Sachschäden in Kauf zu nehmen, wenn dadurch Personenschäden vermeidbar sind.\grqq~\cite{ek}
\end{quote}

\begin{quote}
\glqq
8. Echte dilemmatische Entscheidungen, wie die Entscheidung Leben gegen Leben sind von
der konkreten tatsächlichen Situation unter Einschluss \glqq unberechenbarer\grqq\ Verhaltensweisen Betroffener abhängig. 
Sie sind deshalb nicht eindeutig normierbar und auch nicht
ethisch zweifelsfrei programmierbar. Technische Systeme müssen auf Unfallvermeidung
ausgelegt werden, sind aber auf eine komplexe oder intuitive Unfallfolgenabschätzung
nicht so normierbar, dass sie die Entscheidung eines sittlich urteilsfähigen, verantwortlichen Fahrzeugführers ersetzen 
oder vorwegnehmen könnten. Ein menschlicher Fahrer
würde sich zwar rechtswidrig verhalten, wenn er im Notstand einen Menschen tötet, um
einen oder mehrere andere Menschen zu retten, aber er würde nicht notwendig schuldhaft handeln. Derartige in der Rückschau 
angestellte und besondere Umstände würdigende Urteile des Rechts lassen sich nicht ohne weiteres in abstrakt-generelle 
Ex-AnteBeurteilungen und damit auch nicht in entsprechende Programmierungen umwandeln.
Es wäre gerade deshalb wünschenswert, durch eine unabhängige öffentliche Einrichtung
(etwa einer Bundesstelle für Unfalluntersuchung automatisierter Verkehrssysteme oder
eines Bundesamtes für Sicherheit im automatisierten und vernetzten Verkehr) Erfahrungen systematisch zu verarbeiten.\grqq~\cite{ek}
\end{quote}

\begin{quote}
\glqq
9. Bei unausweichlichen Unfallsituationen ist jede Qualifizierung nach persönlichen Merkmalen (Alter, Geschlecht, 
körperliche oder geistige Konstitution) strikt untersagt. Eine
Aufrechnung von Opfern ist untersagt. Eine allgemeine Programmierung auf eine Minderung der Zahl von Personenschäden 
kann vertretbar sein. Die an der Erzeugung von
Mobilitätsrisiken Beteiligten dürfen Unbeteiligte nicht opfern.\grqq~\cite{ek}
\end{quote}

\begin{quote}
\glqq
10. Die dem Menschen vorbehaltene Verantwortung verschiebt sich bei automatisierten und
vernetzten Fahrsystemen vom Autofahrer auf die Hersteller und Betreiber der technischen Systeme und die infrastrukturellen, 
politischen und rechtlichen Entscheidungsinstanzen. Gesetzliche Haftungsregelungen und ihre Konkretisierung in der gerichtlichen
Entscheidungspraxis müssen diesem Übergang hinreichend Rechnung tragen.\grqq~\cite{ek}
\end{quote}

\begin{quote}
\glqq
11. Für die Haftung für Schäden durch aktivierte automatisierte Fahrsysteme gelten die gleichen Grundsätze wie in der übrigen 
Produkthaftung. Daraus folgt, dass Hersteller oder
Betreiber verpflichtet sind, ihre Systeme fortlaufend zu optimieren und auch bereits ausgelieferte Systeme zu beobachten und zu 
verbessern, wo dies technisch möglich und zumutbar ist.\grqq~\cite{ek}
\end{quote}

\begin{quote}
\glqq
12. Die Öffentlichkeit hat einen Anspruch auf eine hinreichend differenzierte Aufklärung
über neue Technologien und ihren Einsatz. Zur konkreten Umsetzung der hier entwickelten Grundsätze sollten in möglichst 
transparenter Form Leitlinien für den Einsatz und die
Programmierung von automatisierten Fahrzeugen abgeleitet und in der Öffentlichkeit
kommuniziert und von einer fachlich geeigneten, unabhängigen Stelle geprüft werden.\grqq~\cite{ek}
\end{quote}

\begin{quote}
\glqq
13. Ob in Zukunft eine dem Bahn- und Luftverkehr entsprechende vollständige Vernetzung
und zentrale Steuerung sämtlicher Kraftfahrzeuge im Kontext einer digitalen Verkehrsinfrastruktur möglich und sinnvoll sein wird, 
lässt sich heute nicht abschätzen. Eine vollständige Vernetzung und zentrale Steuerung sämtlicher Fahrzeuge im Kontext einer 
digitalen Verkehrsinfrastruktur ist ethisch bedenklich, wenn und soweit sie Risiken einer totalen Überwachung der Verkehrsteilnehmer 
und der Manipulation der Fahrzeugsteuerung nicht sicher auszuschließen vermag.\grqq~\cite{ek}
\end{quote}

\begin{quote}
\glqq
14. Automatisiertes Fahren ist nur in dem Maße vertretbar, in dem denkbare Angriffe, insbesondere Manipulationen des 
IT-Systems oder auch immanente Systemschwächen nicht
zu solchen Schäden führen, die das Vertrauen in den Straßenverkehr nachhaltig erschüttern.\grqq~\cite{ek}
\end{quote}

\begin{quote}
\glqq
15. Erlaubte Geschäftsmodelle, die sich die durch automatisiertes und vernetztes Fahren entstehenden, für die Fahrzeugsteuerung 
erheblichen oder unerheblichen Daten zunutze
machen, finden ihre Grenze in der Autonomie und Datenhoheit der Verkehrsteilnehmer.
Fahrzeughalter oder Fahrzeugnutzer entscheiden grundsätzlich über Weitergabe und
Verwendung ihrer anfallenden Fahrzeugdaten. Die Freiwilligkeit solcher Datenpreisgabe
setzt das Bestehen ernsthafter Alternativen und Praktikabilität voraus. Einer normativen
Kraft des Faktischen, wie sie etwa beim Datenzugriff durch die Betreiber von Suchmaschinen oder sozialen Netzwerken vorherrscht, 
sollte frühzeitig entgegengewirkt werden.\grqq~\cite{ek}
\end{quote}

\begin{quote}
\glqq
16. Es muss klar unterscheidbar sein, ob ein fahrerloses System genutzt wird oder ein Fahrer
mit der Möglichkeit des \glqq Overrulings\grqq\ Verantwortung behält. Bei nicht fahrerlosen Systemen muss die Mensch/Maschine-Schnittstelle 
so ausgelegt werden, dass zu jedem Zeitpunkt klar geregelt und erkennbar ist, welche Zuständigkeiten auf welcher Seite liegen,
insbesondere auf welcher Seite die Kontrolle liegt. Die Verteilung der Zuständigkeiten
(und damit der Verantwortung) zum Beispiel im Hinblick auf Zeitpunkt und Zugriffsregelungen sollte dokumentiert und gespeichert werden. 
Das gilt vor allem für Übergabevorgänge zwischen Mensch und Technik. Eine internationale Standardisierung der Übergabevorgänge 
und der Dokumentation (Protokollierung) ist anzustreben, um angesichts der
grenzüberschreitenden Verbreitung automobiler und digitaler Technologien die Kompatibilität der Protokoll- oder 
Dokumentationspflichten zu gewährleisten.\grqq~\cite{ek}
\end{quote}

\begin{quote}
\glqq
17. Software und Technik hochautomatisierter Fahrzeuge müssen so ausgelegt werden, dass
die Notwendigkeit einer abrupten Übergabe der Kontrolle an den Fahrer (\glqq Notstand\grqq)
praktisch ausgeschlossen ist. Um eine effiziente, zuverlässige und sichere Kommunikation zwischen Mensch und Maschine zu 
ermöglichen und Überforderung zu vermeiden,
müssen sich die Systeme stärker dem Kommunikationsverhalten des Menschen anpassen
und nicht umgekehrt erhöhte Anpassungsleistungen dem Menschen abverlangt werden.\grqq~\cite{ek}
\end{quote}

\begin{quote}
\glqq
18. Lernende und im Fahrzeugbetrieb selbstlernende Systeme sowie ihre Verbindung zu zentralen Szenarien-Datenbanken 
können ethisch erlaubt sein, wenn und soweit sie Sicherheitsgewinne erzielen. Selbstlernende Systeme dürfen nur dann eingesetzt werden, wenn
sie die Sicherheitsanforderungen an fahrzeugsteuerungsrelevante Funktionen erfüllen
und die hier aufgestellten Regeln nicht aushebeln. Es erscheint sinnvoll, relevante Szenarien an einen zentralen 
Szenarien-Katalog einer neutralen Instanz zu übergeben, um
entsprechende allgemeingültige Vorgaben, einschließlich etwaiger Abnahmetests zu erstellen.\grqq~\cite{ek}
\end{quote}

\begin{quote}
\glqq
19. In Notsituationen muss das Fahrzeug autonom, d.h. ohne menschliche Unterstützung, in
einen \glqq sicheren Zustand\grqq\ gelangen. Eine Vereinheitlichung insbesondere der Definition
des sicheren Zustandes oder auch der Übergaberoutinen ist wünschenswert.\grqq~\cite{ek}
\end{quote}

\begin{quote}
\glqq
20. Die sachgerechte Nutzung automatisierter Systeme sollte bereits Teil der allgemeinen digitalen Bildung sein. Der sachgerechte 
Umgang mit automatisierten Fahrsystemen sollte
bei der Fahrausbildung in geeigneter Weise vermittelt und geprüft werden.\grqq~\cite{ek}
\end{quote}

\newpage

\cleardoublepage
\section{Schlussbemerkung}

\newpage

\cleardoublepage
\begin{thebibliography}{99}
\addcontentsline{toc}{section}{Literatur}

\bibitem{ek} Ethik-Kommission, Bundesministerium für Verkehr und digitale Infrastruktur, Automatisiertes und
Vernetztes Fahren, 2017

\bibitem{site1} PRESSEMITTEILUNG - 157/2016 Auftaktsitzung der Ethik-Kommission zum automatisierten Fahren, Stand 30.12.2018, \url{https://www.bmvi.de/SharedDocs/DE/Pressemitteilungen/2016/157-dobrindt-ethikkommission.html}

\end{thebibliography}

\newpage


\end{document}